% FAC-EN-U0052
% French source: fac_body.tex lines 1860--1886.
% Authority: PDF physical pp. 51--52; printed pp. 246--247 (section 2 opening
% and no. 54 complete).
% FAC-FR-C0031 corrects printed "élement" only in corrected French.
% FAC-FR-C0032 corrects the printed f_j/f_i index conflict in corrected
% French and English; diplomatic French preserves the printed j.

\phantomsection\label{fac:III:2}
\section*{\S2. Graded modules and coherent algebraic sheaves on projective space}
\addcontentsline{toc}{section}{Chapter III, \S2. Graded modules and coherent algebraic sheaves on projective space}

\phantomsection\label{fac:III:2:no:54}
\textbf{54. The operation $\mathcal{F}(n)$.}
\phantomsection\label{fac:III:2:no:54:construction:twist}
Let $\mathcal{F}$ be an algebraic sheaf on $X=\mathbf{P}_r(K)$, and let
$\mathcal{F}_i=\mathcal{F}(U_i)$ be the restriction of $\mathcal{F}$ to
$U_i$ (cf. \hyperref[fac:III:1:no:51]{no.~51}). For any integer $n$, let
\phantomsection\label{fac:III:2:no:54:theta}
\[
\theta_{ij}(n):\mathcal{F}_j(U_i\cap U_j)
  \longrightarrow\mathcal{F}_i(U_i\cap U_j)
\]
be the isomorphism defined by multiplication by the function
$t_j^n/t_i^n$. This makes sense because $t_j/t_i$ is a regular function on
$U_i\cap U_j$ with values in $K^*$. At every point of
$U_i\cap U_j\cap U_k$ we have
\[
\theta_{ij}(n)\circ\theta_{jk}(n)=\theta_{ik}(n).
\]
We may therefore apply \hyperref[fac:I:1:no:4:prop:4]{Proposition~4 of no.~4} and obtain an algebraic sheaf,
denoted $\mathcal{F}(n)$, by gluing the sheaves
$\mathcal{F}_i=\mathcal{F}(U_i)$ by means of the isomorphisms
$\theta_{ij}(n)$.

There are canonical isomorphisms
\phantomsection\label{fac:III:2:no:54:formula:canonical-isomorphisms}
\[
\mathcal{F}(0)\simeq\mathcal{F},
\qquad
\mathcal{F}(n)(m)\simeq\mathcal{F}(n+m).
\]
Moreover, $\mathcal{F}(n)$ is locally isomorphic to $\mathcal{F}$, and is
therefore coherent if $\mathcal{F}$ is. It also follows that every exact
sequence
\[
\mathcal{F}\longrightarrow\mathcal{F}'\longrightarrow\mathcal{F}''
\]
of algebraic sheaves gives rise, for every $n\in\mathbf{Z}$, to an exact
sequence
\[
\mathcal{F}(n)\longrightarrow\mathcal{F}'(n)
  \longrightarrow\mathcal{F}''(n).
\]

Applying the \hyperref[fac:III:2:no:54:construction:twist]{preceding construction} to $\mathcal{F}=\mathcal{O}$ gives the
sheaves $\mathcal{O}(n)$, $n\in\mathbf{Z}$. We now give another description
of these sheaves. If $U$ is open in $X$, let $A_U^n$ be the subset of
\[
A_U=\Gamma(\pi^{-1}(U),\mathcal{O}_Y)
\]
formed by the homogeneous functions of degree $n$, that is, the functions
satisfying
\phantomsection\label{fac:III:2:no:54:formula:homogeneous-function}
\[
f(\lambda y)=\lambda^n f(y)
\qquad(\lambda\in K^*,\ y\in\pi^{-1}(U)).
\]
The $A_U^n$ are $A_U^0$-modules and hence give rise to an algebraic sheaf,
which we denote by $\mathcal{O}'(n)$. An element of
$\mathcal{O}'(n)_x$, $x\in X$, may therefore be identified with a rational
fraction $P/Q$, where $P$ and $Q$ are homogeneous polynomials such that
$Q(x)\neq0$ and
\[
\deg P-\deg Q=n.
\]

\phantomsection\label{fac:III:2:prop:1}
\textsc{Proposition 1.}
\textit{The sheaves $\mathcal{O}(n)$ and $\mathcal{O}'(n)$ are canonically
isomorphic.}

By definition, a section of $\mathcal{O}(n)$ over an open subset $U\subset X$
is a system $(f_i)$ of sections of $\mathcal{O}$ over the $U\cap U_i$ such
that
\[
f_i=(t_j^n/t_i^n)f_j
\qquad\text{on }U\cap U_i\cap U_j.
\]
The $f_i$ may be identified with regular homogeneous functions of degree
$0$ on the $\pi^{-1}(U)\cap\pi^{-1}(U_i)$. Put $g_i=t_i^n f_i$. Then
$g_i=g_j$ at every point of
$\pi^{-1}(U)\cap\pi^{-1}(U_i)\cap\pi^{-1}(U_j)$, and hence the $g_i$ are
the restrictions of a unique regular function $g$ on $\pi^{-1}(U)$,
homogeneous of degree $n$. Conversely, such a function $g$ defines a system
$(f_i)$ by $f_i=g/t_i^n$. Thus the map $(f_i)\mapsto g$ is an isomorphism
from $\mathcal{O}(n)$ onto $\mathcal{O}'(n)$.

In what follows we shall most often identify $\mathcal{O}(n)$ and
$\mathcal{O}'(n)$ by means of the \hyperref[fac:III:2:prop:1]{preceding isomorphism}. A section of
$\mathcal{O}'(n)$ over $X$ is simply a regular function on $Y$ homogeneous
of degree $n$. If $r\geq1$, such a function is identically zero for $n<0$,
and is a homogeneous \textit{polynomial} of degree $n$ for $n\geq0$.

\phantomsection\label{fac:III:2:prop:2}
\textsc{Proposition 2.}
\textit{For every algebraic sheaf $\mathcal{F}$, the sheaves
$\mathcal{F}(n)$ and
$\mathcal{F}\otimes_{\mathcal{O}}\mathcal{O}(n)$ are canonically
isomorphic.}

Since $\mathcal{O}(n)$ is obtained from the $\mathcal{O}_i$ by gluing with
the $\theta_{ij}(n)$, the sheaf
$\mathcal{F}\otimes\mathcal{O}(n)$ is obtained from the
$\mathcal{F}_i\otimes\mathcal{O}_i$ by gluing with the isomorphisms
$1\otimes\theta_{ij}(n)$. On identifying
$\mathcal{F}_i\otimes\mathcal{O}_i$ with $\mathcal{F}_i$, we recover exactly
the definition of $\mathcal{F}(n)$.

In what follows we shall likewise identify $\mathcal{F}(n)$ with
$\mathcal{F}\otimes\mathcal{O}(n)$.
